\chapter{Происхождение multiplicity-Hamiltonian из родительского действия}
\label{ch:v8-baryon-c0-multiplicity-environment-hamiltonian-parent-origin-gate}

Предыдущий гейт локализовал минимальный селектор как одноосный Hamiltonian
$h=\varepsilon I_3+\Delta(I_3-P)$. Проверим все уже существующие
квадратичные источники на multiplicity-пространстве.

Старый parent не содержит новых connector-координат, поэтому
\begin{equation}
 H_{\rm old}\big|_{\mathcal E_{\rm mult}}=0_3.
 \label{eq:v8-baryon-c0-parent-origin-old-zero-hamiltonian}
\end{equation}
Единственный след полного $M_5(\mathbb C)$-замыкания даёт только
\begin{equation}
 H_{\tau_5}=\alpha K_{\rm conn}^{(5)}
 =\frac{2\alpha}{5}I_3,
 \qquad P_{\min}(H_{\tau_5})=I_3.
 \label{eq:v8-baryon-c0-parent-origin-m5-scalar-hamiltonian}
\end{equation}
Аналогично полный изотропный процесс имеет
\begin{equation}
 C_{\rm full}=I_3,
 \qquad H_C=\alpha_C I_3,
 \label{eq:v8-baryon-c0-parent-origin-kossakowski-scalar}
\end{equation}
так что ни trace-метрика, ни Kossakowski-ковариация не имеют выделенной
ground line.

\section{Локальные веса допускают, но не выводят анизотропию}

До полного Hom-замыкания концевая алгебра разрешает семейство
\begin{equation}
 H_{\rm loc}(p,q)
 =\operatorname{diag}(p_0+q,p_1+q,p_2+q),
 \qquad p_0+p_1+p_2+2q=1.
 \label{eq:v8-baryon-c0-parent-origin-local-trace-hamiltonian}
\end{equation}
Два положительных нормированных следа
\begin{equation}
 (p_0,p_1,p_2,q)_A=(1/10,1/5,1/2,1/10),
 \qquad
 (p_0,p_1,p_2,q)_B=(1/2,1/10,1/5,1/10)
 \label{eq:v8-baryon-c0-parent-origin-local-trace-witnesses}
\end{equation}
выбирают соответственно разные минимальные coordinate projectors. Значит,
локальная алгебра предоставляет форму допустимого Hamiltonian, но её
трёхмерный trace-симплекс не является селектором.

\section{Почему ранний семейный оператор не переносится автоматически}

На семейном qutrit-носителе Version~4--6 имеется точная матрица
\begin{equation}
 R_4^+=\begin{pmatrix}
 27/2&-2&3/2\\
 -2&35/2&-5/2\\
 3/2&-5/2&17
 \end{pmatrix}.
 \label{eq:v8-baryon-c0-parent-origin-family-r4-plus}
\end{equation}
Она действительно анизотропна и имеет простой спектр, поскольку
\begin{equation}
 \chi_{R_4^+}(\lambda)
 =\frac{4\lambda^3-192\lambda^2+3003\lambda-15358}{4},
 \qquad
 \operatorname{disc}\chi_{R_4^+}=\frac{164241}{16}>0.
 \label{eq:v8-baryon-c0-parent-origin-r4-simple-spectrum}
\end{equation}
Однако текущие multiplicity-координаты являются тремя различными
endpoint-бимодульными метками с проекторами $\Pi_i$. Прямой перенос
нарушает эту диагональность:
\begin{equation}
 \sum_{i=0}^{2}\lVert[R_4^+,\Pi_i]\rVert_{\rm HS}^2=50\ne0.
 \label{eq:v8-baryon-c0-parent-origin-r4-endpoint-commutator}
\end{equation}

Чтобы использовать семейный оператор, необходима типизированная карта
\begin{equation}
 T:\mathcal E_{\rm fam}\longrightarrow\mathcal E_{\rm mult},
 \qquad h_T=T R_4^+T^{\mathsf T}.
 \label{eq:v8-baryon-c0-parent-origin-family-multiplicity-map}
\end{equation}
Но обе трёхмерные multiplicity-реплики тривиальны относительно уже
сопоставленных симметрий, поэтому
\begin{equation}
 \operatorname{Hom}_G(\mathcal E_{\rm fam},\mathcal E_{\rm mult})
 \simeq M_3(\mathbb R),
 \qquad \dim_{\mathbb R}=9,
 \label{eq:v8-baryon-c0-parent-origin-intertwiner-space}
\end{equation}
а после изометрического ограничения остаётся всё $O(3)$. Ни конкретный
$T$, ни его согласование с endpoint-проекторами текущим parent не заданы.
Кроме того, $R_4^+$ без отдельного размерного множителя даёт только
безразмерные относительные щели.

Итоговый аудит источников равен
\begin{equation}
 N_{\rm parent\ origin}=0/5:
 \quad\{H_{\rm old},\tau_5,C_{\rm full},
 \text{local trace simplex},R_4^+\}.
 \label{eq:v8-baryon-c0-parent-origin-candidate-ledger}
\end{equation}

\begin{theorem}[Изотропные источники и нетипизированный near miss]
Все собственные текущие источники на multiplicity-среде равны нулю,
скалярны или принадлежат невыбранному trace-симплексу. Ранний $R_4^+$ имеет
достаточную анизотропию, но действует на другом носителе и не допускает
канонического переноса без новой карты $T$.
\end{theorem}

\begin{nogocriterion}[Запрет неявного отождествления двух qutrit]
Совпадение размерностей $\mathcal E_{\rm fam}$ и $\mathcal E_{\rm mult}$
не задаёт их отождествления. Нельзя использовать $R_4^+$ как Hamiltonian
среды до предъявления gauge-, Real-, grading- и endpoint-совместимого
интертвейнера $T$ и общего размерного коэффициента.
\end{nogocriterion}

\begin{protocol}\begin{enumerate}
\item старый parent на multiplicity-среде: \textbf{$0_3$};
\item полный $M_5$ trace-source: \textbf{скалярный};
\item полный Kossakowski-source: \textbf{скалярный};
\item локальный trace-source: \textbf{анизотропный, но не выбран};
\item $R_4^+$: \textbf{простой спектр, чужой носитель};
\item endpoint-коммутаторный остаток $R_4^+$: \textbf{$50$};
\item пространство формальных межqutrit-карт: \textbf{$M_3(\mathbb R)$};
\item его размерность: \textbf{$9$};
\item выведенных parent-источников: \textbf{$0/5$};
\item направление и абсолютная щель: \textbf{не выведены};
\item следующий гейт: \textbf{допуск family-to-multiplicity intertwiner}.
\end{enumerate}\end{protocol}

Машинный сертификат сохранён в
\path{s2t_v8_baryon_c0_multiplicity_environment_hamiltonian_parent_origin_gate_results.json}.